\documentclass[oneside,a4paper,14pt]{extarticle}
\usepackage[a4paper,letterpaper,top=20mm,bottom=20mm,left=20mm,right=10mm]{geometry}
\usepackage[russian]{babel}
\usepackage{indentfirst}
\usepackage{graphicx}
\usepackage{caption}
\usepackage{titlesec}
\usepackage{minted, fancyvrb}
\usepackage{hyperref}
\usepackage{enumitem}
\usepackage{amsmath}
\usepackage{fancyhdr}
\usepackage{colortbl}
\usepackage[dvipsnames]{xcolor}
\usepackage{float}
\usepackage{tikz-timing}
\usepackage{tikz}
\usetikzlibrary{karnaugh}

% Цвета для карты Карно (чтобы xcolor и tikz их знали)
\definecolor{LogisimKMapColor0}{RGB}{128,0,0}
\definecolor{LogisimKMapColor1}{RGB}{230,25,75}

% Настройка автоматических подписей к рисункам
\DeclareCaptionLabelSeparator{customdash}{~---~}
\captionsetup[figure]{name=Рисунок, labelsep=customdash, justification=centering}

\setminted{style = rainbow_dash, fontsize = \small}

% Форматирование заголовков
\titleformat{\section}{\normalsize\bfseries}{\thesection}{1em}{}
\titleformat{\subsection}{\normalsize\bfseries}{\thesubsection}{1em}{}
\titleformat{\subsubsection}{\normalsize\bfseries}{\thesubsubsection}{1em}{}

% Интерлиньяж и абзац
\renewcommand\baselinestretch{1.33}
\setlength{\parindent}{1.25cm}

% Для всех списков
\setlist[enumerate]{
  left=0.5\parindent,
  label=\arabic*.,
  itemsep=0pt,
  topsep=5pt,
  partopsep=0pt,
  parsep=0pt
}
\setlist[itemize]{
  left=0.5\parindent,
  itemsep=0pt,
  topsep=5pt,
  partopsep=0pt,
  parsep=0pt
}

% Гиперссылки
\hypersetup{
  colorlinks=true,
  linkcolor=black,
  urlcolor=blue,
  pdfborder={0 0 0},
  pdftitle={Контрольная работа: Структурный синтез универсального JK-триггера},
  pdfauthor={Черкасов А.А.},
}

\begin{document}

\newpage
\thispagestyle{empty}
\begin{center}
  МИНИСТЕРСТВО НАУКИ И ВЫСШЕГО ОБРАЗОВАНИЯ РОССИЙСКОЙ ФЕДЕРАЦИИ\\
  ФЕДЕРАЛЬНОЕ ГОСУДАРСТВЕННОЕ БЮДЖЕТНОЕ УЧРЕЖДЕНИЕ ВЫСШЕГО ОБРАЗОВАНИЯ\\
  <<ВЯТСКИЙ ГОСУДАРСТВЕННЫЙ УНИВЕРСИТЕТ>>\\
  Институт математики и информационных систем\\
  Факультет автоматики и вычислительной техники\\
  Кафедра электронных вычислительных машин
\end{center}
\vspace{10mm}

\hfill
\begin{tabular}{l}
  \footnotesize Дата сдачи на проверку:                                          \\
  \footnotesize <<\rule[-1mm]{5mm}{0.10mm}\/>>\rule[-1mm]{20mm}{0.10mm}\ 2026 г. \\
  \footnotesize Проверено:                                                       \\
  \footnotesize <<\rule[-1mm]{5mm}{0.10mm}\/>>\rule[-1mm]{20mm}{0.10mm}\ 2026 г. \\
\end{tabular}
\vfill

\begin{center}
  Отчёт по контрольной работе\\
  по теме\\
  <<Синтез JK-триггера на логических элементах>>\\
  по дисциплине\\
  <<Цифровая схемотехника>>\\
\end{center}
\vspace{25mm}
\noindent
\begin{tabular}{ll}
  Разработал студент гр. ИВТб-2301-05-00 & \hspace{5mm}\rule[-1mm]{30mm}{0.10mm}\,/Черкасов А. А./ \\
                                           & \hspace{12.5mm}\footnotesize(подпись)                   \\
  Преподаватель                            & \hspace{5mm}\rule[-1mm]{30mm}{0.10mm}\,/Мельцов В. Ю./  \\
                                           & \hspace{12.5mm}\footnotesize(подпись)                   \\
\end{tabular}

\noindent
\begin{tabular}{lp{58mm}r}
  Работа защищена &  & \hspace{5mm}<<\rule[-1mm]{5mm}{0.10mm}\/>>\rule[-1mm]{30mm}{0.10mm}\ 2026 г.
\end{tabular}
\vfill

\begin{center}
  Киров\\
  2026
\end{center}

\newpage\thispagestyle{plain}

\section{Цель работы}
Изучить внутреннюю логическую структуру и алгоритмы функционирования бистабильных последовательностных устройств. Выполнить теоретический синтез и моделирование схемы универсального синхронного двухступенчатого JK-триггера со статическими входами асинхронного управления (сброса и установки) на логических вентилях базовых базисов.

\section{Матрица переходов и таблица истинности}
JK-триггер является универсальным автоматом памяти, преобразующим неопределённые состояния RS-типа в счётный инверсный режим переключения при входной комбинации $J=K=1$. Поведение выходной функции триггера в следующий момент времени $Q(t+1)$ определяется текущим состоянием $Q(t)$ и сигналами на информационных входах.

\begin{center}
  Таблица 1 --- Полная таблица переходов JK-триггера \\
  \vspace{3mm}
  \begin{tabular}{ccc|c|c}
    $J$ & $K$ & $Q(t)$ & $Q(t+1)$ & Характер перехода / Режим \\
    \hline
    0 & 0 & 0 & 0 & Режим хранения логического нуля \\
    0 & 0 & 1 & 1 & Режим хранения логической единицы \\
    \hline
    0 & 1 & 0 & 0 & Принудительный сброс в 0 \\
    0 & 1 & 1 & 0 & Принудительный сброс в 0 \\
    \hline
    1 & 0 & 0 & 1 & Принудительная установка в 1 \\
    1 & 0 & 1 & 1 & Принудительная установка в 1 \\
    \hline
    1 & 1 & 0 & 1 & Инверсия состояния (счётный режим) \\
    1 & 1 & 1 & 0 & Инверсия состояния (счётный режим) \\
  \end{tabular}
\end{center}

\section{Минимизация характеристического уравнения триггера}
Для нахождения минимального логического выражения, описывающего логику работы синхронного JK-триггера, построим карту Карно для трёх переменных. Переменными карты являются информационные входы $J, K$ и текущее состояние триггера $Q$.

\begin{center}
  Таблица 2 --- Карта Карно функции переходов $Q(t+1)$ \\
  \vspace{5mm}
  \begin{tikzpicture}[x=1.5cm, y=1.5cm]
    % Сетка карты Карно (2 строки, 4 столбца)
    \draw[thick] (0,0) rectangle (4,2);
    \draw[thick] (1,0) -- (1,2);
    \draw[thick] (2,0) -- (2,2);
    \draw[thick] (3,0) -- (3,2);
    \draw[thick] (0,1) -- (4,1);

    % Разделительная угловая линия для шапки переменных
    \draw[thick] (0,2) -- (-0.6,2.5);
    \node at (-0.15,2.3) {$\quad J, K$};
    \node at (-0.4,1.8) {$Q$};

    % Значения переменных по горизонтали (J, K)
    \node[above] at (0.5,2.05) {00};
    \node[above] at (1.5,2.05) {01};
    \node[above] at (2.5,2.05) {11};
    \node[above] at (3.5,2.05) {10};

    % Значения переменной по вертикали (Q)
    \node[left] at (-0.1,1.5) {0};
    \node[left] at (-0.1,0.5) {1};

    % Заполнение ячеек карты Карно (Минтермы переходов)
    % Строка Q=0
    \node at (0.5,1.5) {0};
    \node at (1.5,1.5) {1};
    \node at (2.5,1.5) {1};
    \node at (3.5,1.5) {1};
    
    % Строка Q=1
    \node at (0.5,0.5) {0};
    \node at (1.5,0.5) {0};
    \node at (2.5,0.5) {0};
    \node at (3.5,0.5) {1};

    % Контуры минимизации (Склеивание единиц)
    % 1. Горизонтальный контур для (J * \bar{Q}) — ячейки 01, 11 при Q=0
    \draw[rounded corners=6pt, fill=LogisimKMapColor1, fill opacity=0.3, draw=LogisimKMapColor1, thick] 
      (1.1,1.1) rectangle (2.9,1.9);

    % 2. Вертикальный контур для (\bar{K} * Q) — крайний правый столбец 10 для Q=0 и Q=1
    \draw[rounded corners=6pt, fill=LogisimKMapColor0, fill opacity=0.3, draw=LogisimKMapColor0, thick] 
      (3.1,0.1) rectangle (3.9,1.9);
  \end{tikzpicture}
\end{center}

В результате выделения и склеивания контуров на карте Карно получаем минимальное характеристическое уравнение переходов JK-триггера в базисе <<И, ИЛИ, НЕ>>:
\begin{equation}
  Q(t+1) = (J \cdot \overline{Q}) \vee (\overline{K} \cdot Q)
\end{equation}

Для реализации полноценного приоритетного управления (входы асинхронного сброса $CLR$ и установки $SET$) и перевода схемы в двухступенчатую структуру (устранение сквозных гонок тактового сигнала $CLK$), в каноническую структуру вводятся дополнительные логические блокировки.

\section{Функциональная схема устройства}
На основании синтезированного уравнения и цепей блокировки была построена и верифицирована функциональная двухступенчатая схема JK-триггера со статическими асинхронными входами установки и сброса.

\begin{figure}[H]
  \centering
  \includegraphics[width=0.95\textwidth,keepaspectratio]{pics/functional_diagram.png}
  \caption{Функциональная схема двухступенчатого JK-триггера}
\end{figure}

\section*{Вывод}
В ходе выполнения контрольной работы были изучены внутренние принципы построения универсальных элементов цифровой памяти. На основе таблицы переходов и аппарата минимизации карт Карно было получено характеристическое уравнение переходов JK-автомата. Спроектированная двухступенчатая функциональная структура полностью подтвердила свою корректность при моделировании, продемонстрировав стабильное переключение без эффекта опасных «гонок» сигналов.

\end{document}